% Part: incompleteness
% Chapter: incompleteness-provability
% Section: rosser-thm

\documentclass[../../../include/open-logic-section]{subfiles}

\begin{document}

\olfileid{inc}{inp}{ros}

\olsection{रोसरचे प्रमेय}

गोडेलची सिद्धता बदलून अधिक सबळ निकाल मिळेल का,
म्हणजे “$\omega$-सुसंगत” ऐवजी फक्त
“सुसंगत” असे गृहीत धरता येईल का?
रोसरने शोधलेल्या युक्तीमुळे याचे उत्तर
“होय” असे आहे. रोसरची युक्ती म्हणजे
$\OProv[\Th{T}](y)$ च्या जागी
“सुधारित” !!{derivability} विधेय
$\ORProv_T(y)$ वापरणे.

\begin{thm}
\ollabel{thm:rosser}
$\Th{T}$ ही~$\Th{Q}$ चा विस्तार
असलेली कोणतीही सुसंगत आणि
!!{axiomatizable} उपपत्ती असू द्या.
मग $\Th{T}$ संपूर्ण नसते.
\end{thm}

\begin{proof}
आठवा, $\OProv[\Th{T}](y)$ ची व्याख्या
$\lexists[x][\OPrf[\Th{T}](x,
  y)]$ अशी आहे. येथे
$\OPrf[\Th{T}](x, y)$ हे अशा निर्णेय
संबंधाचे निरूपण करते की $x$ ही
गोडेल संख्या~$y$ असलेल्या
!!{sentence} च्या !!a{derivation}
ची गोडेल संख्या असते, तर आणि तरच तो
संबंध खरा असतो. $x$ आणि~$y$
यांच्यातील पुढील संबंधही निर्णेय
असतो: $x$ ही गोडेल
संख्या~$y$ असलेल्या वाक्याच्या
\emph{खंडनाची} गोडेल संख्या असते.
$\fn{not}(x)$ हे पुढील कार्य
करणारे आदिम पुनरावर्ती फलन
असू द्या: $x$ हा $!A$ या
सूत्राचा संकेतांक असेल, तर
$\fn{not}(x)$ हा $\lnot
!A$ चा संकेतांक असतो. सूत्रसंकेतांक नसलेल्या
आदानासाठी शून्य परत देऊन हे फलन सर्वत्र परिभाषित करतो. मग
$\Refut[\Th{T}](x, y)$
खरे असते, तर आणि तरच
$\Prf[\Th{T}](x, \fn{not}(y))$
खरे असते. त्याचे निरूपण
$\ORefut[\Th{T}](x, y)$
करू द्या. मग
$\Th{T} \Proves \lnot !A$
असेल आणि $\delta$ ही
संबंधित !!{derivation} असेल,
तर $\Th{Q} \Proves
\ORefut[\Th{T}](\gn{\delta}, \gn{!A})$.
आता $\ORProv[\Th{T}](y)$ ची
व्याख्या अशी करू:
\[
\lexists[x][(\OPrf[\Th{T}](x,y) \land \lforall[z][(z < x \lif \lnot
  \ORefut[\Th{T}](z,y))])].
\]
स्थूलपणे, $\ORProv[\Th{T}](y)$
असे सांगते की “$y$ ची
$\Th{T}$ मध्ये सिद्धता आहे,
आणि कमी संकेतांक असलेले
$y$ चे खंडन नाही.”
$\Th{T}$ सुसंगत असेल,
तर $\ORProv[\Th{T}](y)$ आणि
$\OProv[\Th{T}](y)$ ही दोन्ही
विधाने त्याच संख्यांसाठी
सत्य असतात. पण $\Th{T}$
मधील \emph{सिद्धतायोग्यतेच्या}
दृष्टीने (आता सत्य आणि
सिद्धतायोग्यता यांत फरक आहे
हे आपल्याला माहीत आहे)
त्यांचे गुणधर्म भिन्न असतात.
$\Th{T}$ \emph{वि}संगत
असेल, तर ती दोन्ही विधाने
त्याच संख्यांसाठी सत्य
\emph{नाहीत}!{}
($\ORProv[\Th{T}](y)$
याचा अर्थ अनेकदा “$y$ रोसर
अर्थाने सिद्धतायोग्य आहे”
असा घेतला जातो. पण आत्ताच
पाहिल्याप्रमाणे रोसर
सिद्धतायोग्यता हा
सिद्धतायोग्यतेचा काही
विशेष प्रकार नाही:
विसंगत उपपत्तीत काही
!!{sentence}s सिद्धतायोग्य
असतात, पण रोसर अर्थाने
सिद्धतायोग्य नसतात.
म्हणून हा शब्द गोंधळात
टाकू शकतो. तो टाळण्यासाठी
“$y$ शमोव्हेबल आहे”
असेही म्हणता येईल.)

स्थिरबिंदू उपपत्तीवरून
असे सूत्र $!R_\Th{T}$
असते की
\begin{equation}
  \Th{Q} \Proves !R_\Th{T} \liff \lnot \ORProv[\Th{T}](\gn{!R_\Th{T}}).
  \ollabel{RT}
\end{equation}
\olref[1in]{thm:first-incompleteness}
च्या सिद्धतेपेक्षा इथे
आपला दावा अधिक सबळ आहे:
$\Th{T}$ सुसंगत असेल,
तर $\Th{T}$ हे
$!R_\Th{T}$ ला !!{derive}
करत नाही, आणि $\Th{T}$
$\lnot !R_\Th{T}$ लाही
!!{derive} करत नाही.
(म्हणजे $\omega$-सुसंगततेचे
गृहीतक लागत नाही.)

प्रथम $\Th{T} \Proves/
!R_{\Th{T}}$ हे दाखवू.
उलट ते सिद्ध होत असेल,
तर~$T$ पासून
$!R_{\Th{T}}$ ची
!!a{derivation} आहे;
तिची गोडेल संख्या $n$
मानू. $\OPrf[\Th{T}]$
हे $\Prf[\Th{T}]$ चे
$\Th{Q}$ मध्ये निरूपण
करत असल्यामुळे
$\Th{Q} \Proves
\OPrf[\Th{T}](\num{n}, \gn{!R_{\Th{T}}})$.
शिवाय प्रत्येक $k < n$
साठी, $\Th{T}$ सुसंगत
असल्यामुळे, $k$ ही
$\lnot !R_{\Th{T}}$
च्या !!a{derivation} ची
गोडेल संख्या नसते.
म्हणून प्रत्येक $k < n$
साठी $\Th{Q} \Proves \lnot
\ORefut[\Th{T}](\num{k}, \gn{!R_{\Th{T}}})$.
\olref[req][min]{lem:less-nsucc}
वरून $\Th{Q} \Proves
\lforall[z][(z < \num{n} \lif \lnot \ORefut[\Th{T}](z,
  \gn{!R_{\Th{T}}}))]$.
त्यामुळे
\[
\Th{Q} \Proves \lexists[x][(\OPrf[\Th{T}](x,\gn{!R_{\Th{T}}}) \land \lforall[z][(z
    < x \lif \lnot \ORefut[\Th{T}](z,\gn{!R_{\Th{T}}}))])],
\]
पण हेच
$\ORProv[\Th{T}](\gn{!R_{\Th{T}}})$
आहे. \olref{RT} वरून
$\Th{Q}
\Proves \lnot !R_{\Th{T}}$.
$\Th{T}$ हा $\Th{Q}$ चा
विस्तार असल्यामुळे
$\Th{T}
\Proves \lnot !R_{\Th{T}}$
सुद्धा मिळते. पण आपण
$\Th{T} \Proves !R_{\Th{T}}$
असे गृहीत धरले होते.
त्यामुळे $\Th{T}$ विसंगत
ठरेल; हे प्रमेयाच्या
गृहीतकाच्या विरुद्ध आहे.

आता $\Th{T} \Proves/
\lnot !R_{\Th{T}}$ हे
दाखवू. उलट ते सिद्ध होत
असेल आणि
$\lnot !R_{\Th{T}}$
च्या !!a{derivation} ची
गोडेल संख्या $n$ असेल
असे मानू. मग
$\Refut[\Th{T}](n, \Gn{!R_{\Th{T}}})$
खरे असते. $\ORefut[\Th{T}]$
हे $\Refut[\Th{T}]$ चे
$\Th{Q}$ मध्ये निरूपण
करत असल्यामुळे
$\Th{Q} \Proves
\ORefut[\Th{T}](\num{n}, \gn{!R_{\Th{T}}})$.
मग $\Th{T}$ हे
$!R_{\Th{T}}$ सुद्धा
!!{derive} करेल आणि
त्यामुळे विसंगत ठरेल,
हे पुन्हा दाखवू. कारण
\begin{align*}
\Th{Q} & \Proves !R_{\Th{T}} \liff \lnot \ORProv[\Th{T}](\gn{!R_{\Th{T}}}),
\intertext{आणि $\Th{T}$ हा~$\Th{Q}$ चा विस्तार असल्यामुळे पुढील दाखवणे पुरेसे आहे:}
\Th{Q} & \Proves \lnot \ORProv[\Th{T}](\gn{!R_{\Th{T}}}).
\end{align*}
!!{sentence} $\lnot
\ORProv[\Th{T}](\gn{!R_{\Th{T}}})$,
म्हणजे
\begin{align*}
  \lnot & \lexists[x][(\OPrf[\Th{T}](x,\gn{!R_{\Th{T}}}) \land
    \lforall[z][(z < x \lif \lnot \ORefut[\Th{T}](z,\gn{!R_{\Th{T}}}))])],
  \intertext{हे तर्कशास्त्रीय दृष्ट्या पुढील विधानाशी सममूल्य आहे:}
  & \lforall[x][(\OPrf[\Th{T}](x,\gn{!R_{\Th{T}}}) \lif
    \lexists[z][(z < x \land \ORefut[\Th{T}](z,\gn{!R_{\Th{T}}}))])].
\end{align*}
$\Th{Q}$ काय !!{derive}s
करते याविषयीची तथ्ये वापरून
अनौपचारिक युक्तिवाद करू.
$x$ स्वेच्छेने निवडलेला
आणि $\OPrf[\Th{T}](x,
\gn{!R_{\Th{T}}})$ खरा
आहे असे मानू. आपल्याला
आधीच माहीत आहे की
$\Th{T} \Proves/ !R_{\Th{T}}$.
म्हणून प्रत्येक $k$
साठी $\Th{Q} \Proves \lnot
\OPrf[\Th{T}](\num{k}, \gn{!R_{\Th{T}}})$.
त्यामुळे प्रत्येक $k$
साठी $\eq/[x][\num{k}]$
मिळते. विशेषतः (a)
$\eq/[x][\num{n}]$
मिळते. तसेच
$\lnot(\eq[x][\num{0}]
\lor \eq[x][\num{1}] \lor \dots
\lor \eq[x][\num{n-1}])$
मिळते. म्हणून
\olref[req][min]{lem:less-nsucc}
वरून (b) $\lnot(x < \num{n})$.
\olref[req][min]{lem:trichotomy}
वरून $\num{n} < x$.
$\Th{Q} \Proves
\ORefut[\Th{T}](\num{n}, \gn{!R_{\Th{T}}})$
असल्यामुळे
$\num{n} < x \land
\ORefut[\Th{T}](\num{n}, \gn{!R_{\Th{T}}})$
मिळते, आणि त्यावरून
$\lexists[z][(z < x
\land \ORefut[\Th{T}](z,\gn{!R_{\Th{T}}}))]$
मिळते. $x$ स्वेच्छेने
निवडलेला असल्यामुळे
अपेक्षेप्रमाणे पुढील मिळते:
\[
\lforall[x][(\OPrf[\Th{T}](x,\gn{!R_{\Th{T}}}) \lif
  \lexists[z][(z < x \land \ORefut[\Th{T}](z,\gn{!R_{\Th{T}}}))])].
\]
\end{proof}

\begin{prob}
नैसर्गिक संख्यांचे दोन संच
$A$ आणि $B$ यांना
\emph{संगणनक्षमपणे अविलगनीय}
म्हणतात, जर $A \subseteq X$
आणि $B \subseteq \Complement{X}$
असेल असा कोणताही
!!{decidable} संच~$X$
नसेल ($\Complement{X}$
म्हणजे~$X$ चा पूरक
$\Nat \setminus X$).
$\Th{T}$ ही~$\Th{Q}$
ची सुसंगत,
!!{axiomatizable} विस्तार-उपपत्ती
असू द्या. $\Th{T}$
मध्ये सिद्ध होणाऱ्या
!!{sentence}s च्या गोडेल
संख्यांचा संच $A$,
आणि $\Th{T}$ मध्ये
खंडित होणाऱ्या वाक्यांच्या
गोडेल संख्यांचा संच $B$
असेल असे मानू. $A$ आणि
$B$ हे संगणनक्षमपणे
अविलगनीय आहेत हे सिद्ध करा.
\end{prob}

\end{document}
